\section{Vocabulaire}

    \subsection{Expressions}

        \subsubsection{Les salutations (\ruby{挨拶}{あいさつ})}

    \begin{longtblr}{%
        colspec = {Q[l,2] Q[l,3]},  
        width = \linewidth,
        hlines = {0.6pt, gray!50},
        vlines = {0pt},
        row{1} = {font=\bfseries, bg=couleurmatiere!20}
    }
        日本語 & Signification \\

        お{早}{はよ}う(ございます) & bonmatin \\
        \ruby{今日は}{こんにちわ} & bonjour \\
        \ruby{今晩は}{こんばんわ} & bonsoir \\
        お\ruby{休み}{やすみ}(なさい) & bonne nuit \\

        さようなら & au revoir \\
        またね & à bientôt, à plus \\
        また\ruby{明日}{あした} & à demain \\
        また\ruby{今度}{こんど} & à la prochaine \\
        じゃあ、また & à plus, à bientôt \\

        \ruby{初}{はじ}めまして & enchanté \\
        (どうぞ)よろしくお\ruby{願}{ねが}いします & ravi de vous rencontrer \\

        いらっしゃいませ & bienvenue (dans notre magasin) \\

        ありがとう(ございます) & merci \\
        どういたしまして & de rien \\
        すみません & excusez-moi, pardon \\
        ごめん(なさい) & désolé, pardon \\

        いただきます & bon appétit / formule avant de manger \\
        ごちそうさまでした & merci pour le repas / formule après avoir mangé \\
    \end{longtblr}

        \subsubsection{\ruby{紹介}{しょうかい}する}

        \subsubsection{À trier}

    \begin{tblr}[theme=general]{colspec = {Q[l,2] Q[l,4] Q[l,6]}}
        漢字 & 仮名 & Signification \\

        & もういちど & encore une fois \\
        & あのお & hmm... \\
        & そうですね & exactement, c’est ça \\
        & そうですか & je vois... vraiment ? \\
        & じゃあ & d’accord... si c’est comme ça... \\
        & どうぞ & voici, je vous en prie \\
        & どうも & merci \\
        (\~を) お願いします & (\~を) おねがいします & s’il vous plaît (pas de réponse attendue) \\
        (\~を) ください & s’il vous plaît (réponse attendue) \\
    \end{tblr}


    \subsection{Noms par par thématique}

        \subsubsection{L’entourage}

    \begin{longtblr}{%
        colspec = {Q[l,1] Q[l,2]},  
        width = \linewidth,
        hlines = {0.6pt, gray!50},
        vlines = {0pt},
        row{1} = {font=\bfseries, bg=couleurmatiere!20}
    }
        日本語 & Personne \\

        \ruby{両親}{りょうしん} & parents \\
        \ruby{母}{はは} & (ma) maman \\
        お\ruby{母}{かあ}さん & (votre) mère \\
        \ruby{父}{ちち} & (mon) papa \\
        お\ruby{父}{とう}さん & (votre) père \\

        \ruby{妻}{つま} & (ma) femme \\
        \ruby{奥}{おく}さん & (votre) femme \\
        \ruby{夫}{おっと} & (mon) mari \\
        ご\ruby{主人}{しゅじん} & (votre) mari \\

        \ruby{兄弟}{きょうだい} & sœurs et frères \\
        \ruby{姉}{あね} & (ma) sœur aînée \\
        お\ruby{姉}{ねえ}さん & (votre) sœur aînée \\
        \ruby{妹}{いもうと} & (ma) sœur cadette \\
        \ruby{妹}{いもうと}さん & (votre) sœur cadette \\
        \ruby{兄}{あに} & (mon) frère aîné \\
        お\ruby{兄}{にい}さん & (votre) frère aîné \\
        \ruby{弟}{おとうと} & (mon) frère cadet \\
        \ruby{弟}{おとうと}さん & (votre) frère cadet \\

        \ruby{子供}{こども} & enfant \\
        \ruby{娘}{むすめ} & (ma) fille \\
        \ruby{娘}{むすめ}さん & (votre) fille \\
        \ruby{息子}{むすこ} & (mon) fils \\
        \ruby{息子}{むすこ}さん & (votre) fils \\

        \ruby{祖母}{そぼ} & (ma) grand-mère \\
        お\ruby{祖母}{ばあ}さん & (votre) grand-mère \\
        \ruby{祖父}{そふ} & (mon) grand-père \\
        お\ruby{祖父}{じい}さん & (votre) grand-père \\

        \ruby{友達}{ともだち} & amie \\
        \ruby{先生}{せんせい} & professeur / enseignant \\
    \end{longtblr}

        \subsubsection{Les lieux}

    \begin{longtblr}{%
        colspec = {Q[l,1] Q[l,2]},  
        width = \linewidth,
        hlines = {0.6pt, gray!50},
        vlines = {0pt},
        row{1} = {font=\bfseries, bg=couleurmatiere!20}
    }
        日本語 & Lieu \\

        \ruby{市役所}{しやくしょ} & mairie \\
        \ruby{区役所}{くやくしょ} & mairie d'arrondissement \\ 
        \ruby{学校}{がっこう} & école \\ 
        \ruby{小学校}{しょうがっこう} & école primaire \\ 
        \ruby{中学校}{ちゅうがっこう} & collège \\ 
        \ruby{高校}{こうこう} & lycée \\ 
        \ruby{大学}{だいがく} & université \\ 
        \ruby{図書館}{としょかん} & bibliothèque \\ 
        \ruby{本屋}{ほんや} & librairie ; maison d'édition \\ 
        \ruby{美術館}{びじゅつかん} & musée (art) \\ 
        \ruby{博物館}{はくぶつかん} & musée \\ 
        \ruby{映画館}{えいがかん} & cinéma \\ 
        \ruby{劇場}{げきじょう} & théâtre / salle de spectacle \\ 
        コンサートホール & salle de concert \\ 
        \ruby{銀行}{ぎんこう} & banque \\ 
        \ruby{郵便局}{ゆうびんきょく} & bureau de poste \\ 
        \ruby{病院}{びょういん} & hôpital \\ 
        \ruby{診療所}{しんりょうじょ} & cabinet médical / clinique \\ 
        \ruby{薬局}{やっきょく} & pharmacie \\ 
        \ruby{駅}{えき} & gare \\ 
        \ruby{空港}{くうこう} & aéroport 
        \\ バス\ruby{停}{てい} & arrêt de bus \\ 
        \ruby{公園}{こうえん} & parc \\ 
        \ruby{庭}{にわ} & jardin \\ 
        \ruby{広場}{ひろば} & place publique \\ 
        \ruby{駐車場}{ちゅうしゃじょう} & parking \\ 
        \ruby{道}{みち} & chemin, voie, route \\ 
        \ruby{市場}{いちば} & marché \\ 
        スーパー & supermarché (supermarket) \\ 
        コンビニ & supérette (convenience store) \\ 
        デパート & grand magasin (department store) \\ 
        レストラン & restaurant \\ 
        ホテル & hôtel \\ 
        \ruby{神社}{じんじゃ} & sanctuaire shinto \\ 
        お\ruby{寺}{てら} & temple bouddhiste\\ 
        \ruby{教会}{きょうかい} & église \\ 
        \ruby{体育館}{たいいくかん} & gymnase \\ 
        プール & piscine \\ 
        \ruby{運動場}{うんどうじょう} & terrain de sport\\
    \end{longtblr}

        \subsubsection{Les couleurs}
    
    \begin{longtblr}{%
        colspec = {Q[l,1] Q[l,2]},  
        width = \linewidth,
        hlines = {0.6pt, gray!50},
        vlines = {0pt},
        row{1} = {font=\bfseries, bg=couleurmatiere!20}
    }
        日本語 & Signification \\

        \ruby{色}{いろ} & couleur \\
        \ruby{赤}{あか} & rouge \\
        \ruby{青}{あお} & bleu \\
        \ruby{黄色}{きいろ} & jaune \\
        \ruby{緑}{みどり} & vert \\
        \ruby{紫}{むらさき} & violet, pourpre \\
        \ruby{白}{しろ} & blanc \\
        \ruby{黒}{くろ} & noir \\
        \ruby{灰色}{はいいろ} & gris \\
        \ruby{茶色}{ちゃいろ} & marron / brun \\
        \ruby{桃色}{ももいろ} & rose \\
        \ruby{橙色}{だいだいいろ} & orange \\
        \ruby{金色}{きんいろ} & doré, or \\
        \ruby{銀色}{ぎんいろ} & argenté, argent \\
    \end{longtblr}

        \subsubsection{Les professions et statuts}

    \begin{longtblr}{%
        colspec = {Q[l,1] Q[l,2]},  
        width = \linewidth,
        hlines = {0.6pt, gray!50},
        vlines = {0pt},
        row{1} = {font=\bfseries, bg=couleurmatiere!20}
    }
        日本語 & Signification \\

        \ruby{仕事}{しごと} & travail / métier \\
        \ruby{会社}{かいしゃ} & compagnie, entreprise, société \\
        \ruby{会社員}{かいしゃいん} & salarié / employé \\
        \ruby{学生}{がくせい} & étudiant \\
        \ruby{大学生}{だいがくせい} & étudiant universitaire \\
        \ruby{高校生}{こうこうせい} & lycéen \\
        \ruby{中学生}{ちゅうがくせい} & collégien \\
        \ruby{小学生}{しょうがくせい} & écolier / élève du primaire \\
        \ruby{教師}{きょうし} & enseignant / professeur \\
        \ruby{医者}{いしゃ} & médecin \\
        \ruby{看護師}{かんごし} & infirmier / infirmière \\
        \ruby{薬剤師}{やくざいし} & pharmacien / pharmacienne \\
        \ruby{警察官}{けいさつかん} & policier / policière \\
        \ruby{運転手}{うんてんしゅ} & chauffeur / conducteur \\
        \ruby{店員}{てんいん} & employé de magasin / vendeur \\
        \ruby{店長}{てんちょう} & gérant / responsable de magasin \\
        \ruby{銀行員}{ぎんこういん} & employé de banque \\
        \ruby{公務員}{こうむいん} & fonctionnaire \\
        \ruby{社長}{しゃちょう} & président / patron d'entreprise \\
        \ruby{料理人}{りょうりにん} & cuisinier \\
        \ruby{美容師}{びようし} & coiffeur / coiffeuse \\
        \ruby{農家}{のうか} & agriculteur / agricultrice \\
        \ruby{歌手}{かしゅ} & chanteur / chanteuse \\
        \ruby{俳優}{はいゆう} & acteur / actrice \\
        \ruby{作家}{さっか} & écrivain / auteur \\
        \ruby{記者}{きしゃ} & journaliste / reporter \\
        エンジニア & ingénieur \\
    \end{longtblr}

        \subsubsection{Les animaux}

    \begin{longtblr}{%
        colspec = {Q[l,1] Q[l,2]},  
        width = \linewidth,
        hlines = {0.6pt, gray!50},
        vlines = {0pt},
        row{1} = {font=\bfseries, bg=couleurmatiere!20}
    }
        日本語 & Signification \\

        \ruby{動物}{どうぶつ} & animal \\ 
        \ruby{犬}{いぬ} & chien \\ 
        \ruby{猫}{ねこ} & chat \\ 
        \ruby{鳥}{とり} & oiseau \\ 
        \ruby{魚}{さかな} & poisson \\ 
        \ruby{馬}{うま} & cheval \\ 
        \ruby{牛}{うし} & vache, bœuf \\ 
        \ruby{豚}{ぶた} & cochon \\ 
        \ruby{羊}{ひつじ} & mouton \\ 
        \ruby{山羊}{やぎ} & chèvre \\ 
        \ruby{兎}{うさぎ} & lapin \\ 
        \ruby{鼠}{ねずみ} & souris / rat \\ 
        \ruby{猿}{さる} & singe \\ 
        \ruby{狐}{きつね} & renard \\ 
        \ruby{狸}{たぬき} & tanuki \\ 
        \ruby{鶏}{にわとり} & poule, poulet \\ 
        \ruby{鴨}{かも} & canard \\ 
        \ruby{蜂}{はち} & abeille / guêpe \\
    \end{longtblr}

        \subsubsection{Éléments naturels, météo et saisons}

    \begin{longtblr}{%
        colspec = {Q[l,1] Q[l,2]},  
        width = \linewidth,
        hlines = {0.6pt, gray!50},
        vlines = {0pt},
        row{1} = {font=\bfseries, bg=couleurmatiere!20}
    }
        日本語 & Signification \\

        \ruby{水}{みず} & eau \\ 
        \ruby{火}{ひ} & feu \\ 
        \ruby{木}{き} & arbre, bois \\ 
        \ruby{土}{つち}、\ruby{地}{ち} & terre, sol \\ 
        \ruby{風}{かぜ} & vent \\ 
        \ruby{雨}{あめ} & pluie \\ 
        \ruby{雪}{ゆき} & neige \\ 
        \ruby{氷}{こおり} & glace \\ 
        \ruby{空気}{くうき} & air \\ 
        \ruby{空}{そら} & ciel \\ 
        \ruby{世界}{せかい} & monde, univers \\ 
        \ruby{地球}{ちきゅう} & Terre \\ 
        \ruby{海}{うみ} & mer \\ 
        \ruby{川}{かわ} & rivière / fleuve \\ 
        \ruby{田}{た} & rizière \\ 
        \ruby{山}{やま} & montagne \\ 
        \ruby{坂}{さか} & colline, pente \\ 
        \ruby{石}{いし} & pierre \\ 
        \ruby{砂}{すな} & sable \\ 
        \ruby{草}{くさ} & herbe \\ 
        \ruby{雲}{くも} & nuage \\ 
        \ruby{太陽}{たいよう} & soleil \\ 
        \ruby{月}{つき} & lune \\ 
        \ruby{星}{ほし} & étoile \\ 
        \ruby{光}{ひかり} & lumière \\ 
        \ruby{影}{かげ} & ombre \\ 
        \ruby{雷}{かみなり} & tonnerre / foudre \\ 
        \ruby{台風}{たいふう} & typhon \\ 
        \ruby{嵐}{あらし} & tempête \\ 
        \ruby{霧}{きり} & brouillard \\ 
        \ruby{波}{なみ} & vague \\ 
        \ruby{温度}{おんど} & température \\

        \ruby{季節}{きせつ} & saison \\
        \ruby{春}{はる} & printemps \\
        \ruby{夏}{なつ} & été \\
        \ruby{秋}{あき} & automne \\
        \ruby{冬}{ふゆ} & hiver \\
    \end{longtblr}

        \subsubsection{Moyens de transport}

    \begin{longtblr}{%
        colspec = {Q[l,1] Q[l,2]},  
        width = \linewidth,
        hlines = {0.6pt, gray!50},
        vlines = {0pt},
        row{1} = {font=\bfseries, bg=couleurmatiere!20}
    }
        日本語 & Signification \\

        \ruby{交通}{こうつう} & transport / circulation \\
        \ruby{電車}{でんしゃ} & train électrique \\
        \ruby{新幹線}{しんかんせん} & Shinkansen, train à grande vitesse \\
        \ruby{地下鉄}{ちかてつ} & métro \\
        \ruby{鉄道}{てつどう} & chemin de fer, réseau ferroviaire \\
        \ruby{駅}{えき} & gare \\
        \ruby{車}{くるま} & voiture \\
        \ruby{自動車}{じどうしゃ} & automobile, voiture \\
        \ruby{駐車場}{ちゅうしゃじょう} & parking \\
        \ruby{道路}{どうろ} & route \\
        \ruby{高速道路}{こうそくどうろ} & autoroute \\
        タクシー & taxi \\
        バス & bus \\
        バス\ruby{停}{てい} & arrêt de bus \\
        \ruby{自転車}{じてんしゃ} & vélo \\
        バイク & moto \\
        \ruby{飛行機}{ひこうき} & avion \\
        \ruby{空港}{くうこう} & aéroport \\
        \ruby{船}{ふね} & bateau, navire \\
    \end{longtblr}

        \subsubsection{Parties du corps}

    \begin{longtblr}{%
        colspec = {Q[l,1] Q[l,2]},  
        width = \linewidth,
        hlines = {0.6pt, gray!50},
        vlines = {0pt},
        row{1} = {font=\bfseries, bg=couleurmatiere!20}
    }
        日本語 & Signification \\

        \ruby{体}{からだ} & corps \\
        \ruby{頭}{あたま} & tête \\
        \ruby{顔}{かお} & visage \\
        \ruby{髪}{かみ} & cheveux \\
        \ruby{目}{め} & œil \\
        \ruby{耳}{みみ} & oreille \\
        \ruby{鼻}{はな} & nez \\
        \ruby{口}{くち} & bouche \\
        \ruby{歯}{は} & dent \\
        \ruby{舌}{した} & langue \\
        \ruby{声}{こえ} & voix \\
        \ruby{首}{くび} & cou \\
        \ruby{肩}{かた} & épaule \\
        \ruby{腕}{うで} & bras \\
        \ruby{手}{て} & main \\
        \ruby{指}{ゆび} & doigt \\
        \ruby{爪}{つめ} & ongle \\
        \ruby{胸}{むね} & poitrine / torse \\
        \ruby{背中}{せなか} & dos \\
        お\ruby{腹}{なか}& ventre \\
        \ruby{腰}{こし} & taille, bas du dos \\
        \ruby{足}{あし} & jambe, pied \\
        \ruby{膝}{ひざ} & genou \\
        \ruby{足首}{あしくび} & cheville \\
        \ruby{踵}{かかと} & talon \\

        \ruby{心臓}{しんぞう} & cœur \\

    \end{longtblr}

        \subsubsection{Habits}

    \begin{longtblr}{%
        colspec = {Q[l,2] Q[l,4] Q[l,6]},  
        width = \linewidth,
        hlines = {0.6pt, gray!50},
        vlines = {0pt},
        row{1} = {font=\bfseries, bg=couleurmatiere!20}
    }
        漢字 & 仮名 & Signification \\

        服 & ふく & vêtement / habits \\
        洋服 & ようふく & vêtements de style occidental \\
        着物 & きもの & kimono \\
        浴衣 & ゆかた & yukata \\
        & シャツ & chemise / T-shirt \\
        & Tシャツ & T-shirt \\
        & セーター & pull / sweater \\
        & コート & manteau \\
        & ジャケット & veste / blouson \\
        & ズボン & pantalon \\
        & ジーンズ & jean \\
        & スカート & jupe \\
        & ドレス & robe \\
        下着 & したぎ & sous-vêtements \\
        靴 & くつ & chaussures \\
        靴下 & くつした & chaussettes \\
        & スニーカー & baskets / chaussures de sport \\
        帽子 & ぼうし & chapeau / casquette \\
        手袋 & てぶくろ & gants \\
        & スカーフ & écharpe \\
        & ネクタイ & cravate \\
        & ベルト & ceinture \\
        & めがね & lunettes \\
        水着 & みずぎ & maillot de bain \\
        & パジャマ & pyjama \\
        鞄 & かばん & sac / cartable \\
        傘 & かさ & parapluie \\
    \end{longtblr}

        \subsubsection{La maison}

    \begin{longtblr}{%
        colspec = {Q[l,2] Q[l,4] Q[l,6]},  
        width = \linewidth,
        hlines = {0.6pt, gray!50},
        vlines = {0pt},
        row{1} = {font=\bfseries, bg=couleurmatiere!20}
    }
        漢字 & 仮名 & Signification \\

        家 & いえ / うち & maison (bâtiment / foyer) \\
        家具 & かぐ & mobilier / meubles \\
        机 & つくえ & bureau / table de travail \\
        椅子 & いす & chaise \\
        & テーブル & table \\
        & ベッド & lit \\
        布団 & ふとん & futon / literie \\
        & ソファ & canapé \\
        棚 & たな & étagère / meuble de rangement \\
        本棚 & ほんだな & bibliothèque / étagère à livres \\
        箪笥 & たんす & commode / armoire à tiroirs \\
        鏡 & かがみ & miroir \\
        電気 & でんき & lumière / électricité \\
        電灯 & でんとう & lampe / éclairage électrique \\
        時計 & とけい & horloge \\
        冷蔵庫 & れいぞうこ & réfrigérateur \\
        電子レンジ & でんしレンジ & four à micro-ondes \\
        洗濯機 & せんたくき & machine à laver \\
        掃除機 & そうじき & aspirateur \\
        & テレビ & télévision \\
        エアコン & エアコン & climatisation / chauffage (air conditionner) \\
        扇風機 & せんぷうき & ventilateur \\
        & ドア & porte \\
        鍵 & かぎ & clé / serrure \\
        窓 & まど & fenêtre \\
        カーテン & カーテン & rideau (curtain) \\
        & トイレ & toilettes \\
    \end{longtblr}

        \subsubsection{Les objets du quotidien}

    \begin{longtblr}{%
        colspec = {Q[l,2] Q[l,4] Q[l,6]},  
        width = \linewidth,
        hlines = {0.6pt, gray!50},
        vlines = {0pt},
        row{1} = {font=\bfseries, bg=couleurmatiere!20}
    }
        漢字 & 仮名 & Signification \\

        電話 & でんわ & téléphone \\
        携帯電話 & けいたいでんわ & téléphone portable \\
        & パソコン & ordinateur (personal computer)\\
        & カメラ & appareil photo \\
        & ラジオ & radio \\
        財布 & さいふ & portefeuille / porte-monnaie \\
        本 & ほん & livre \\
        新聞 & しんぶん & journal \\
        雑誌 & ざっし & magazine / revue \\
        写真 & しゃしん & photo / photographie \\
        箱 & はこ & boîte / caisse \\
        紙 & かみ & papier \\
        鉛筆 & えんぴつ & crayon à papier \\
        & ペン & stylo \\
        消しゴム & けしゴム & gomme \\
    \end{longtblr}

        \subsubsection{La nourriture}

    \begin{longtblr}{%
        colspec = {Q[l,2] Q[l,4] Q[l,6]},  
        width = \linewidth,
        hlines = {0.6pt, gray!50},
        vlines = {0pt},
        row{1} = {font=\bfseries, bg=couleurmatiere!20}
    }
        漢字 & 仮名 & Signification \\

        食べ物 & たべもの & nourriture / aliment \\
        飲み物 & のみもの & boisson \\
        料理 & りょうり & plat, cuisine \\
        ご飯 & ごはん & riz (cuit), repas \\
        朝ご飯 & あさごはん & petit déjeuner \\
        昼ご飯 & ひるごはん & déjeuner \\
        晩ご飯 & ばんごはん & dîner \\
        米 & こめ & riz (cru) \\
        & パン & pain \\
        肉 & にく & viande \\
        牛肉 & ぎゅうにく & bœuf \\
        豚肉 & ぶたにく & porc \\
        鶏肉 & とりにく & poulet / viande de poulet \\
        魚 & さかな & poisson \\
        卵 & たまご & œuf \\
        野菜 & やさい & légumes \\
        果物 & くだもの & fruits \\
        リンゴ & りんご & pomme \\
        & バナナ & banane \\
        苺 & いちご & fraise \\
        牛乳 & ぎゅうにゅう & lait \\
        お茶 & おちゃ & thé \\
        紅茶 & こうちゃ & thé noir \\
        緑茶 & りょくちゃ & thé vert (japonais) \\
        & コーヒー & café \\
        & ジュース & jus / boisson fruitée \\
        お酒 & おさけ & alcool / saké \\
        & ビール & bière \\
        塩 & しお & sel \\
        砂糖 & さとう & sucre \\
        醤油 & しょうゆ & sauce soja \\
        味噌 & みそ & miso \\
        寿司 & すし & sushi \\
    \end{longtblr}

        \subsubsection{Les arts}

    \begin{longtblr}{%
        colspec = {Q[l,2] Q[l,4] Q[l,6]},  
        width = \linewidth,
        hlines = {0.6pt, gray!50},
        vlines = {0pt},
        row{1} = {font=\bfseries, bg=couleurmatiere!20}
    }
        漢字 & 仮名 & Signification \\

        芸術 & げいじゅつ & art \\
        美術 & びじゅつ & beaux-arts / art \\
        絵 & え & dessin / peinture / image \\
        絵画 & かいが & peinture / œuvre picturale \\
        絵本 & えほん & livre d'images / album illustré \\
        写真 & しゃしん & photographie \\
        彫刻 & ちょうこく & sculpture \\
        作品 & さくひん & œuvre \\
        画家 & がか & peintre \\
        芸術家 & げいじゅつか & artiste \\
        音楽 & おんがく & musique \\
        歌 & うた & chanson / chant \\
        歌手 & かしゅ & chanteur / chanteuse \\
        楽器 & がっき & instrument de musique \\
        動画 & どうが & vidéo, animation \\
        映画 & えいが & film / cinéma \\
        映画館 & えいがかん & cinéma (salle) \\
        劇 & げき & pièce de théâtre / spectacle \\
        劇場 & げきじょう & théâtre / salle de spectacle \\
        演劇 & えんげき & théâtre (art dramatique) \\
        俳優 & はいゆう & acteur / actrice \\
        小説 & しょうせつ & roman \\
        文学 & ぶんがく & littérature \\
        詩 & し & poésie / poème \\
        漫画 & まんが & manga / bande dessinée \\
        美術館 & びじゅつかん & musée des beaux-arts \\
        博物館 & はくぶつかん & musée \\
        展覧会 & てんらんかい & exposition \\
        展示 & てんじ & exposition / présentation d'œuvres \\
        文化 & ぶんか & culture \\
        伝統 & でんとう & tradition \\
    \end{longtblr}

        \subsubsection{Les sciences}

    \begin{longtblr}{%
        colspec = {Q[l,2] Q[l,4] Q[l,6]},  
        width = \linewidth,
        hlines = {0.6pt, gray!50},
        vlines = {0pt},
        row{1} = {font=\bfseries, bg=couleurmatiere!20}
    }
        漢字 & 仮名 & Signification \\

        科学 & かがく & science \\
        理科 & りか & sciences naturelles \\
        生物 & せいぶつ & biologie / être vivant \\
        植物 & しょくぶつ & plante / végétal \\
        動物 & どうぶつ & animal \\
        実験 & じっけん & expérience / expérimentation \\
        研究 & けんきゅう & recherche / étude \\
        研究室 & けんきゅうしつ & laboratoire / salle de recherche \\
        実験室 & じっけんしつ & laboratoire \\
        観察 & かんさつ & observation \\
        調査 & ちょうさ & enquête / étude / investigation \\
        結果 & けっか & résultat \\
        方法 & ほうほう & méthode / manière \\
        理由 & りゆう & raison \\
        資料 & しりょう & document / donnée / matériel de référence \\
        & データ & données \\
        道具 & どうぐ & outil / matériel \\
        機械 & きかい & machine \\
        器具 & きぐ & instrument / équipement \\
    \end{longtblr}

        \subsubsection{Autres}

    \begin{tblr}[theme=general]{colspec = {Q[l,2] Q[l,4] Q[l,6]}}
        漢字 & 仮名 & Signification \\

        一度 & いちど & une fois, une opportunité \\
        一番 & いちばん & le plus; premier, meilleur \\
        一部 & いちぶ & une partie, une portion \\

        友達 & ともだち & ami \\
        出身 & しゅっしん & origine, lieu de naissance \\
        皆さん & みなさん & tout le monde \\
    \end{tblr}

    \subsection{Adjectifs par thématiques} 

        \subsubsection{Caractéristiques physiques}

    \begin{tblr}[theme=general]{colspec = {Q[l,2] Q[l,4] Q[l,6]}}
        漢字 & 仮名 & Signification \\

        大きい & おおきい & grand / gros \\
        小さい & ちいさい & petit \\
        高い & たかい & grand (en taille), haut \\
        低い & ひくい & petit (en taille), bas \\
        長い & ながい & long \\
        短い & みじかい & court \\
        太い & ふとい & gros, épais \\
        細い & ほそい & mince, fin \\
        重い & おもい & lourd \\
        軽い & かるい & léger \\
        広い & ひろい & large, vaste \\
        狭い & せまい & étroit \\
        強い & つよい & fort, solide \\
        弱い & よわい & faible, fragile \\

        \ruby{明}{あかる}い & éclairé, clair, éblouissant \\
    \end{tblr}

        \subsubsection{Émotions}

    \begin{longtblr}{%
        colspec = {Q[l,2] Q[l,4] Q[l,6]},  
        width = \linewidth,
        hlines = {0.6pt, gray!50},
        vlines = {0pt},
        row{1} = {font=\bfseries, bg=couleurmatiere!20}
    }
        日本語 & Signification \\

        \ruby{楽}{たの}しい & joyeux, amusant \\
        \ruby{嬉}{うれ}しい & heureux, content \\
        \ruby{明}{あか}るい & clair, gai, radieux \\
        \ruby{悲}{かな}しい & triste \\
        \ruby{寂}{さび}しい & seul, triste de solitude \\

        \ruby{怖}{こわ}い & effrayant, avoir peur \\
        \ruby{恐}{おそ}ろしい & terrifiant, épouvantable \\
        \ruby{面白}{おもしろ}い & intéressant, amusant \\
        \ruby{詰}{つま}らない & ennuyeux, sans intérêt \\
        \ruby{恥}{は}ずかしい & gêné, honteux \\
        \ruby{悔}{くや}しい & frustré, contrarié \\
        \ruby{苦}{くる}しい & pénible, douloureux \\
        \ruby{痛}{いた}い & douloureux, avoir mal \\
        \ruby{眠}{ねむ}い & avoir sommeil \\
        \ruby{忙}{いそが}しい & occupé, débordé \\
        \ruby{疲}{つか}れた & fatigué \\
        \ruby{元気}{げんき} & en forme, énergique \\
        \ruby{静}{しず}か & calme, tranquille \\
        \ruby{穏}{おだ}やか & calme, paisible \\
        \ruby{幸}{しあわ}せ & heureux, épanoui \\
        \ruby{不幸}{ふこう} & malheureux \\
        \ruby{心配}{しんぱい} & inquiet, préoccupé \\
        \ruby{安心}{あんしん} & rassuré, en sécurité \\
        \ruby{不安}{ふあん} & anxieux, inquiet \\
        \ruby{大丈夫}{だいじょうぶ} & aller bien, être en bon état \\
        \ruby{暇}{ひま} & libre, ne pas être occupé \\
    \end{longtblr}

        \subsubsection{Temps et espace}

    \begin{tblr}[theme=general]{colspec = {Q[l,2] Q[l,4] Q[l,6]}}
        漢字 & 仮名 & Signification \\

        早い & はやい & tôt, rapide (temps) \\
        速い & はやい & rapide (vitesse) \\
        遅い & おそい & tard, lent \\
        長い & ながい & long \\
        短い & みじかい & court \\
        近い & ちかい & proche, près \\
        遠い & とおい & loin, éloigné, distant \\
        新しい & あたらしい & nouveau, récent \\
        古い & ふるい & ancien / vieux \\
        若い & わかい & jeune \\
        有り勝ち (な) & ありがち & fréquent, commun, usuel
        珍しい & めずらしい & rare, inhabituel \\
    \end{tblr}

        \subsubsection{}

    \begin{tblr}[theme=general]{colspec = {Q[l,2] Q[l,4] Q[l,6]}}
        漢字 & 仮名 & Signification \\

        美味しい & おいしい & délicieux, bon, savoureux
    \end{tblr}

    

    \begin{tblr}[theme=general]{colspec = {Q[l,2] Q[l,4] Q[l,6]}}
        漢字 & 仮名 & Signification \\

        
    \end{tblr}

    \subsection{Les adverbes}

        \subsubsection{Temps, fréquence et espace}

    \begin{tblr}[theme=general]{colspec = {Q[l,2] Q[l,4] Q[l,6]}}
        漢字 & 仮名 & Signification \\

        & いつも & toujours, tout le temps \\
        大抵 & たいてい & habituellement, généralement \\
        良く & よく & souvent \\
        時々 & ときどき & de temps en temps \\
        一度も & いちども & jamais (いちど + particule d’inclusion も) \\



        & ゆっくり & lentement, doucement \\
    \end{tblr}

        \subsubsection{Quantité}

    \begin{longtblr}{%
        colspec = {Q[l,2] Q[l,4] Q[l,6]},  
        width = \linewidth,
        hlines = {0.6pt, gray!50},
        vlines = {0pt},
        row{1} = {font=\bfseries, bg=couleurmatiere!20}
    }
        漢字 & 仮名 & Signification \\



        沢山 & たくさん & beaucoup \\
        いっぱい & いっぱい & beaucoup / plein \\
        少し & すこし & un peu \\
        一寸 & ちょっと & un peu \\
        全然 & ぜんぜん & pas du tout \\
        全部 & ぜんぶ & tout, entièrement \\
        皆 & みんな & tout le monde, tous \\
        & ほとんど & presque, la plupart \\

        & すごく & extrêmement \\
        十分 & じゅうぶん & suffisamment / assez \\
        & もっと & davantage / plus \\
        & あまり & peu \\
        & ずっと & beaucoup plus / nettement plus \\
        少なくとも & すくなくとも & au moins \\
        多く & おおく & beaucoup / en grand nombre \\
        大体 & だいたい & environ / à peu près \\
        約 & やく & environ \\
        半分 & はんぶん & moitié \\
    \end{longtblr}

        \subsubsection{Manière}

    

        \subsubsection{À trier}

    \begin{tblr}[theme=general]{colspec = {Q[l,2] Q[l,4] Q[l,6]}}
        漢字 & 仮名 & Signification \\

        
        良く & よく & souvent, bien \\

        一緒 & いっしょ & ensemble \\
        各各 & それぞれ & chacun, respectivement \\
        

    \end{tblr}

    \subsection{Les verbes}

        \subsubsection{Verbes sino-japonais}

    \begin{tblr}[theme=general]{colspec = {Q[l,2] Q[l,4] Q[l,6]}}
        日本語 & Signification \\

        \ruby{質問}{しつもん}する & interroger, questionner \\
        \ruby{理解}{りかい}する & comprendre \\
        \ruby{挨拶}{あいさつ}する & saluer, dire bonjour \\
        
    \end{tblr}

        \subsubsection{Quotidien}

    \begin{longtblr}{%
        colspec = {Q[l,1] Q[l,2]},  
        width = \linewidth,
        hlines = {0.6pt, gray!50},
        vlines = {0pt},
        row{1} = {font=\bfseries, bg=couleurmatiere!20}
    }
        日本語 & Signification \\

        \ruby{起}{お}きる & se lever, se réveiller \\
        \ruby{休}{やす}む & se reposer, prendre congé \\
        \ruby{寝}{ね}る & dormir, se coucher \\

        \ruby{着}{き}る & mettre, porter (vêtement du haut) \\
        \ruby{履}{は}く & mettre, porter (chaussures, pantalon) \\
        \ruby{脱}{ぬ}ぐ & enlever (un vêtement) \\
        

        \ruby{食}{た}べる & manger \\
        \ruby{飲}{の}む & boire \\

        \ruby{洗}{あら}う & laver \\
        \ruby{掃除}{そうじ}する & nettoyer \\
        \ruby{洗濯}{せんたく}する & faire la lessive \\
        \ruby{料理}{りょうり}する & cuisiner \\

        \ruby{働}{はたら}く & travailler \\
        \ruby{勉強}{べんきょう}する & étudier \\
        \ruby{練習}{れんしゅう}する & s'entraîner, pratiquer \\

        \ruby{使}{つか}う & utiliser \\
        \ruby{作}{つく}る & faire, fabriquer, préparer \\
        \ruby{開}{あ}ける & ouvrir \\
        \ruby{開}{あ}く & s'ouvrir \\
        \ruby{閉}{し}める & fermer \\
        \ruby{閉}{し}まる & se fermer \\
        \ruby{置}{お}く & poser, mettre \\
        \ruby{持}{も}つ & tenir, avoir, posséder \\
        \ruby{取}{と}る & prendre \\
    \end{longtblr}

        \subsubsection{Déplacement}

    \begin{longtblr}{%
        colspec = {Q[l,1] Q[l,2]},  
        width = \linewidth,
        hlines = {0.6pt, gray!50},
        vlines = {0pt},
        row{1} = {font=\bfseries, bg=couleurmatiere!20}
    }
        日本語 & Signification \\

        \ruby{行}{い}く & aller \\
        \ruby{来}{く}る & venir \\
        \ruby{帰}{かえ}る & rentrer, retourner \\
        \ruby{歩}{ある}く & marcher \\
        \ruby{走}{はし}る & courir \\
        \ruby{乗}{の}る & monter dans, prendre (un transport) \\
        \ruby{降}{お}りる & descendre (d'un transport) \\
        \ruby{渡}{わた}る & traverser \\
        \ruby{曲}{ま}がる & tourner \\
        \ruby{止}{と}まる & s'arrêter \\
        \ruby{出}{で}る & sortir, partir (intransitif) \\
        \ruby{出}{だ}す & sortir, mettre dehors (transitif) \\
        \ruby{入}{はい}る & entrer \\
        \ruby{着}{つ}く & arriver \\
        \ruby{住}{す}む & habiter \\
    \end{longtblr}

        \subsubsection{Interactions sociales et communications}

    \begin{longtblr}{%
        colspec = {Q[l,1] Q[l,2]},  
        width = \linewidth,
        hlines = {0.6pt, gray!50},
        vlines = {0pt},
        row{1} = {font=\bfseries, bg=couleurmatiere!20}
    }
        日本語 & Signification \\

        \ruby{会}{あ}う & rencontrer \\
        \ruby{聞}{き}く & écouter, demander \\
        \ruby{呼}{よ}ぶ & appeler, inviter \\
        \ruby{待}{ま}つ & attendre \\
        \ruby{助}{たす}ける & aider \\
        \ruby{手伝}{てつだ}う & aider, donner un coup de main \\
        \ruby{貸}{か}す & prêter \\
        \ruby{借}{か}りる & emprunter \\
        \ruby{送}{おく}る & envoyer, accompagner \\
        もらう & recevoir \\
        あげる & donner \\
        \ruby{見}{み}せる & montrer \\

        \ruby{話}{はな}す & parler \\
        \ruby{言}{い}う & dire \\
        \ruby{教}{おし}える & enseigner, informer \\
        \ruby{習}{なら}う & apprendre (auprès de quelqu'un) \\
        \ruby{伝}{つた}える & transmettre, faire savoir \\
        \ruby{質問}{しつもん}する & poser une question \\
        \ruby{答}{こた}える & répondre \\
        \ruby{分}{わ}かる & comprendre \\
        \ruby{知}{し}る & savoir, connaître \\
        \ruby{思}{おも}う & penser, trouver que \\
        \ruby{忘}{わす}れる & oublier \\
        \ruby{覚}{おぼ}える & mémoriser, retenir \\
        \ruby{考}{かんが}える & réfléchir, penser \\
    \end{longtblr}

        \subsubsection{Loisirs et sports}

    \begin{longtblr}{%
        colspec = {Q[l,1] Q[l,2]},  
        width = \linewidth,
        hlines = {0.6pt, gray!50},
        vlines = {0pt},
        row{1} = {font=\bfseries, bg=couleurmatiere!20}
    }
        日本語 & Signification \\

        \ruby{遊}{あそ}ぶ & jouer, s'amuser \\
        \ruby{楽}{たの}しむ & profiter de, s'amuser \\
        \ruby{旅行}{りょこう}する & voyager \\
        \ruby{見物}{けんぶつ}する & visiter, faire du tourisme \\
        \ruby{見}{み}る & regarder, voir \\
        \ruby{聞}{き}く & écouter \\
        \ruby{読}{よ}む & lire \\
        \ruby{書}{か}く & écrire \\
        \ruby{歌}{うた}う & chanter \\
        \ruby{撮}{と}る & prendre (une photo) \\
        \ruby{踊}{おど}る & danser \\
    



        \ruby{泳}{およ}ぐ & nager \\
        \ruby{走}{はし}る & courir \\
        \ruby{登}{のぼ}る & monter, escalader \\
    \end{longtblr}

        \subsubsection{Courses et achats}

    \begin{longtblr}{%
        colspec = {Q[l,1] Q[l,2]},  
        width = \linewidth,
        hlines = {0.6pt, gray!50},
        vlines = {0pt},
        row{1} = {font=\bfseries, bg=couleurmatiere!20}
    }
        日本語 & Signification \\

        \ruby{買}{か}う & acheter \\
        \ruby{売}{う}る & vendre \\
        \ruby{払}{はら}う & payer \\
        \ruby{選}{えら}ぶ & choisir \\
        \ruby{注文}{ちゅうもん}する & commander \\
        \ruby{食}{た}べる & manger \\
        \ruby{飲}{の}む & boire \\
        \ruby{頼}{たの}む & commander, demander \\
        \ruby{持}{も}つ & avoir, tenir \\
        \ruby{要}{い}る & avoir besoin de \\
        \ruby{借}{か}りる & emprunter, louer \\
        \ruby{貸}{か}す & prêter, louer \\
        \ruby{返}{かえ}す & rendre, restituer \\
    \end{longtblr}
      

    


        \subsubsection{À trier}

    \begin{tblr}[theme=general]{colspec = {Q[l,1] Q[l,2]}}
        日本語 & Signification \\

        \ruby{使}{つか}う & employer, utiliser \\
        \ruby{住}{す}む & habiter, vivre, résider \\
        \ruby{知}{し}る & savoir, connaître \\
    \end{tblr}


    \subsection{Autres}

        \subsubsection{Pronoms personnels}

    \begin{longtblr}{%
        colspec = {Q[l,2] Q[l,4] Q[l,6]},  
        width = \linewidth,
        hlines = {0.6pt, gray!50},
        vlines = {0pt},
        row{1} = {font=\bfseries, bg=couleurmatiere!20}
    }
        日本語 & Signification \\

        \ruby{私}{わたし} & je, moi \\
        \ruby{僕}{ぼく} & je, moi (masculin, courant) \\
        \ruby{俺}{おれ} & je, moi (masculin, familier) \\
        うち & je moi (féminin) \\

        あなた & tu, vous \\
        \ruby{君}{きみ} & tu (familier, souvent masculin) \\

        \ruby{彼}{かれ} & il, lui \\
        \ruby{彼女}{かのじょ} & elle \\

        \ruby{私}{わたし}たち & nous \\
        \ruby{僕}{ぼく}たち & nous (masculin) \\
        \ruby{俺}{おれ}たち & nous (masculin, familier) \\

        あなたたち & vous (pluriel) \\

        \ruby{彼}{かれ}ら & ils, eux \\
        \ruby{彼女}{かのじょ}たち & elles, elles \\

        \ruby{誰}{だれ} & qui \\

        \ruby{皆}{みな} & tout le monde, tous \\
        \ruby{皆}{みんな} & tout le monde (courant) \\
    \end{longtblr}

    

        \subsubsection{Mots interrogatifs}

    \begin{longtblr}{%
        colspec = {Q[l,2] Q[l,4] Q[l,6]},  
        width = \linewidth,
        hlines = {0.6pt, gray!50},
        vlines = {0pt},
        row{1} = {font=\bfseries, bg=couleurmatiere!20}
    }
        漢字 & 仮名 & Signification \\

        何 & なに / なん & quoi, quel \\ 
        誰 & だれ & qui \\ 
        誰の & だれの & à qui, de qui \\ 
        & どこ & où \\ 
        & どちら & lequel, où (poli) \\ 
        & どっち & lequel, où (familier) \\ 
        & いつ & quand \\ 
        & なぜ & pourquoi \\ 
        & どうして & pourquoi, comment se fait-il que \\ 
        & どう & comment, de quelle manière \\ 
        & どの & quel (devant un nom) \\ 
        & どんな & quel genre de \\ 
        & どれ & lequel \\ 
        & どのくらい & combien, dans quelle mesure \\ 
        & いくら & combien (prix) \\ 
        & いくつ & combien, quel âge \\ 
        何歳 & なんさい & quel âge \\ 
        何時 & なんじ & quelle heure \\ 
        何分 & なんぷん & combien de minutes \\ 
        何曜日 & なんようび & quel jour de la semaine \\ 
        何月 & なんがつ & quel mois \\ 
        何日 & なんにち & quel jour, combien de jours \\ 
        何人 & なんにん & combien de personnes \\ 
        何回 & なんかい & combien de fois \\ 
        & どうやって & comment, par quel moyen \\ 
        どのくらいの時間 & どのくらいのじかん & combien de temps \\
    \end{longtblr}

        \subsubsection{Position dans l’espace}

    \begin{longtblr}{%
        colspec = {Q[l,2] Q[l,4] Q[l,6]},  
        width = \linewidth,
        hlines = {0.6pt, gray!50},
        vlines = {0pt},
        row{1} = {font=\bfseries, bg=couleurmatiere!20}
    }
        漢字 & 仮名 & Signification \\

        上 & うえ & dessus, au-dessus \\
        下 & した & dessous, en dessous \\
        前 & まえ & devant, avant \\
        後ろ & うしろ & derrière \\
        中 & なか & dans, à l'intérieur \\
        外 & そと & dehors, à l'extérieur \\
        右 & みぎ & droite \\
        左 & ひだり & gauche \\
        隣 & となり & à côté, voisin \\
        近く & ちかく & près, à proximité \\
        遠く & とおく & loin, au loin \\
        間 & あいだ & entre \\
    \end{longtblr}

        \subsubsection{La météo}
        
    \begin{longtblr}{%
        colspec = {Q[l,2] Q[l,4] Q[l,6]},  
        width = \linewidth,
        hlines = {0.6pt, gray!50},
        vlines = {0pt},
        row{1} = {font=\bfseries, bg=couleurmatiere!20}
    }
        漢字 & 仮名 & Signification \\

        天気 & てんき & météo, temps \\ 
        天気予報 & てんきよほう & prévisions météo \\ 
        天気がいい & てんきがいい & il fait beau \\
        天気が悪い & てんきがわるい & il fait mauvais \\
        晴れる & はれる & se dégager, faire beau \\ 
        曇る & くもる & se couvrir \\ 
        雨 & あめ & pluie \\ 
        雨が降る & あめがふる & pleuvoir \\ 
        雪 & ゆき & neige \\ 
        雪が降る & ゆきがふる & neiger \\ 
        雪が積もる & ゆきがつもる & s’accumuler (la neige) \\
        凍る & こおる & geler \\
        風 & かぜ & vent \\ 
        風が吹く & かぜがふく & souffler (pour le vent) \\ 
        強風 & きょうふう & vent fort \\ 
        台風 & たいふう & typhon \\ 
        雷 & かみなり & tonnerre, éclair \\ 
        気温 & きおん & température (de l'air) \\ 
        暑い & あつい & chaud \\ 
        暖かい & あたたかい & doux, chaud \\ 
        涼しい & すずしい & frais \\ 
        寒い & さむい & froid \\ 
        雲 & くも & nuage \\ 
        太陽 & たいよう & soleil \\ 
        虹 & にじ & arc-en-ciel \\ 
        霧 & きり & brouillard \\ 
        暑くなる & あつくなる & se réchauffer \\ 
        寒くなる & さむくなる & se refroidir \\ 
    \end{longtblr}

        \subsubsection{Mots de coordination}

    \begin{tblr}[theme=general]{colspec = {Q[l,2] Q[l,4] Q[l,6]}}
        漢字 & 仮名 & Signification \\

        & そして & ensuite, puis, et \\
        & でも & mais, cependant; non plus \\
    \end{tblr}


    \subsection{Les nombres}

    \begin{tblr}[theme = general]{colspec = {Q[c,2] Q[l,2] Q[l,2] Q[l,2]}} 
        数字 & 漢字 & 音読み & 訓読み \\
        0 & 零 & レイ / ゼロ &  \\
        1 & 一 & \underline{イチ} & ひと \\ 
        2 & 二 & \underline{ニ} & ふた \\
        3 & 三 & \underline{サン} & み \\
        4 & 四 & シ & \underline{よん} / よ \\
        5 & 五 & \underline{ゴ} & いつ \\
        6 & 六 & \underline{ロク} & む \\
        7 & 七 & シチ & \underline{なな} \\
        8 & 八 & \underline{ハチ} & や \\
        9 & 九 & ク / \underline{キュウ} & ここの \\
        10 & 十 & \underline{ジュウ} & とお \\
        100 & 百 & \underline{ヒャク} &  \\
        1000 & 千 & \underline{セン} &  \\
    \end{tblr}

    \begin{tblr}[theme = general]{colspec = {Q[c,2] Q[l,2] Q[l,2] Q[l,2]}} 
        数字 & 漢字 & 音読み \\
        $10^4$ & 万  & マン \\
        $10^8$ & 億 & オク \\
        $10^{12}$ & 兆 & チョウ \\
    \end{tblr}

        \subsubsection{Compter}

    On peut à l'aide du premier tableau de kanjis et prononciations (dont on préférera celles soulignées lorsque l'on compte) déjà compter jusqu'à 9999.

    \begin{tcolorbox}
        -- 七千四百六十五

        -- ななセン　よんヒャク　ロクジュウ　ゴ

        -- 7465
    \end{tcolorbox}

    Quelques changements de prononciation sont dus à l'assemblage des kanjis, comme 
    
    \begin{tblr}[theme = general]{colspec = {Q[c,2] Q[l,2] Q[l,2]}} 
        数字 & 漢字 & 読み \\
        300 & 三百 & サン\textcolor{couleurmatiere}{ビ}ャク \\
        600 & 六白 & ロ\textcolor{couleurmatiere}{ッピ}ャク \\
        800 & 八白 & ハ\textcolor{couleurmatiere}{ッピ}ャク \\
        3000 & 三千 & サン\textcolor{couleurmatiere}{ゼ}ン \\
        8000 & 八千 & ハ\textcolor{couleurmatiere}{ッ}セン \\
    \end{tblr}

    Les puissances par multiple de 4 de 10 permettent ensuite de former des paquets avec lesquels on pourra compter jusqu'à l'infini. 

    \begin{tcolorbox}
        -- 二千三百四十五\textcolor{couleurmatiere}{兆}六千七百八十九\textcolor{couleurmatiere}{億}百二十三\textcolor{couleurmatiere}{万}四千五百六十七

        -- ニセンサンビャクよんジュウゴ　チョウ　ロクセンななヒャクハチジュウキュウ　オク　ヒャクニジュウサン　マン　よんセンゴヒャクロクジュウなな

        -- 2345 6789 0123 4567
    \end{tcolorbox}

    \subsection{La date}

    Pour décrire le temps, on utilisera majoritairement 6 kanjis : 年 (ネン, année), 月 (ガツ, mois), 日 (ニチ ou か ou び, jour), 時 (ジ, heure), 分 (フン, minutes), 秒 (ビョウ, seconde).

        \subsubsection{L'heure}

    À la façon de l'anglais, on distingue les heures avant midi (午前, ごぜん) de celles après midi (午後, ごご). 

    \begin{longtblr}{%
        colspec = {Q[l,2] Q[l,4] Q[l,6]},  
        width = \linewidth,
        hlines = {0.6pt, gray!50},
        vlines = {0pt},
        row{1} = {font=\bfseries, bg=couleurmatiere!20}
    }
        Heure & 漢字 & 仮名 \\
        1h & 午前一時 & ごぜんイチジ \\
        2h & 午前二時 & ごぜんニジ \\
        3h & 午前三時 & ごぜんサンジ \\
        4h & 午前四時 & ごぜん\textcolor{couleurmatiere}{よ}ジ \\
        5h & 午前五時 & ごぜんゴジ \\
        6h & 午前六時 & ごぜんロクジ \\
        7h & 午前七時 & ごぜん\textcolor{couleurmatiere}{シチ}ジ \\
        8h & 午前八時 & ごぜんハチジ \\
        9h & 午前九時 & ごぜん\textcolor{couleurmatiere}{ク}ジ \\
        10h & 午前十時 & ごぜんジュウジ \\
        11h & 午前十一時 & ごぜんジュウイチジ \\
        12h & 十二時 & ジュウニジ \\
        13h & 午後一時 & ごごイチジ \\
        14h & 午後二時 & ごごニジ \\
        15h & 午後三時 & ごごサンジ \\
        16h & 午後四時 & ごご\textcolor{couleurmatiere}{よ}ジ \\
        17h & 午後五時 & ごごゴジ \\
        18h & 午後六時 & ごごロクジ \\
        19h & 午後七時 & ごご\textcolor{couleurmatiere}{シチ}ジ \\
        20h & 午後八時 & ごごハチジ \\
        21h & 午後九時 & ごご\textcolor{couleurmatiere}{ク}ジ \\
        22h & 午後十時 & ごごジュウジ \\
        23h & 午後十一時 & ごごジュウイチジ \\
        0h & 零時 & レイジ \\
    \end{longtblr}

    Pour les minutes, on utilisera les mêmes prononciation des nombres que pour compter, avec une prise d'accent ou non à la prononciation de フン :

     \begin{tblr}[theme = general]{colspec = {Q[c,2] Q[l,2] Q[l,3]}} 
        Minute & 漢字 & 仮名 \\
        1min & 一分 & イ\textcolor{couleurmatiere}{ップン} \\
        2min & 二分 & ニフン \\
        3min & 三分 & サン\textcolor{couleurmatiere}{プン} \\
        4min & 四分 & よん\textcolor{couleurmatiere}{プン} \\
        5min & 五分 & ゴフン \\
        6min & 六分 & ロ\textcolor{couleurmatiere}{ップン} \\
        7min & 七分 & ななフン \\
        8min & 八分 & ハ\textcolor{couleurmatiere}{ップン} \\
        9min & 九分 & キュウフン \\
        10min & 十分 & ジュ\textcolor{couleurmatiere}{ップン} \\
    \end{tblr}

    En ce qui concerne les secondes, on utilisera les mêmes prononciations des nombres que pour compter, auxquelles on ajoute ビョウ.

    \begin{tcolorbox}
        -- 午前七時四十六分五十三秒

        -- 午前７時４６分５３秒

        -- ゴゼンシチジ　よんジュウロップン　ゴジュウサンビョウ

        -- 7h46min53s
    \end{tcolorbox}

        \subsubsection{Les jours}

    Les jours de la semaine sont formés à partir du nom des astres (Soleil, Lune, Mars...) auxquels on additionne le nom « jour de la semaine » (曜日, ようび).

    \begin{tblr}[theme = general]{colspec = {Q[c,2] Q[l,2] Q[l,3]}} 
        Jour & 漢字 & 仮名 \\
        Dimanche & 日曜日 & にちようび \\
        Lundi & 月曜日 & げつようび \\
        Mardi & 火曜日 & かようび \\
        Mercredi & 水曜日 & すいようび \\
        Jeudi & 木曜日 & もくようび \\
        Vendredi & 金曜日 & きんようび \\
        Samedi & 土曜日 & どようび \\
    \end{tblr}

    Pour ce qui est du numéro d'un jour, on utilisera un décompte particulier.

    \begin{longtblr}{colspec = {Q[c,2] Q[l,2] Q[l,3]}, width = \linewidth,
        hlines = {0.6pt, gray!50},
        vlines = {0pt},
        row{1} = {font=\bfseries, bg=couleurmatiere!20}} 
        Mois & 漢字 & 仮名 \\
        1er & 一日 / 1日 & \textcolor{couleurmatiere}{ついたち} (nouvelle lune) \\
        2 & 二日 / 2日 & \textcolor{couleurmatiere}{ふつ}\textcolor{gray}{か} \\
        3 & 三日 / 3日 & \textcolor{couleurmatiere}{みっ}\textcolor{gray}{か} \\
        4 & 四日 / 4日 & \textcolor{couleurmatiere}{よっ}\textcolor{gray}{か} \\
        5 & 五日 / 5日 & \textcolor{couleurmatiere}{いつ}\textcolor{gray}{か} \\
        6 & 六日 / 6日 & \textcolor{couleurmatiere}{むい}\textcolor{gray}{か} \\
        7 & 七日 / 7日 & \textcolor{couleurmatiere}{なの}\textcolor{gray}{か} \\
        8 & 八日 / 8日 & \textcolor{couleurmatiere}{よ}\textcolor{gray}{か} \\
        9 & 九日 / 9日 & \textcolor{couleurmatiere}{ここの}\textcolor{gray}{か} \\
        10 & 十日 / 10日 & \textcolor{couleurmatiere}{とう}\textcolor{gray}{か} \\
        11 & 十一日 / 11日 & ジュウイチニチ \\
        12 & 十二日 / 12日 & ジュウニニチ \\
        13 & 十三日 / 13日 & ジュウサンニチ \\
        14 & 十四日 / 14日 & ジュウ\textcolor{couleurmatiere}{よっ}\textcolor{gray}{か} \\
        15 & 十五日 / 15日 & ジュウゴニチ \\
        16 & 十六日 / 16日 & ジュウロクニチ \\
        17 & 十七日 / 17日 & ジュウシチニチ \\
        18 & 十八日 / 18日 & ジュウハチニチ \\
        19 & 十九日 / 19日 & ジュウ\textcolor{couleurmatiere}{キュウ}ニチ \\
        20 & 二十日 / 20日 & はつ\textcolor{gray}{か} \\
        21 & 二十一日 / 21日 & ニジュウイチニチ \\
        22 & 二十二日 / 22日 & ニジュウニニチ \\
        23 & 二十三日 / 23日 & ニジュウサンニチ \\
        24 & 二十四日 / 24日 & ニジュウ\textcolor{couleurmatiere}{よっ}\textcolor{gray}{か} \\
        25 & 二十五日 / 25日 & ニジュウゴニチ \\
        26 & 二十六日 / 26日 & ニジュウロクニチ \\
        27 & 二十七日 / 27日 & ニジュウシチニチ \\
        28 & 二十八日 / 28日 & ニジュウハチニチ \\
        29 & 二十九日 / 29日 & ニジュウ\textcolor{couleurmatiere}{キュウ}ニチ \\
        30 & 三十日 / 30日 & サンジュウニチ \\
        31 & 三十一日 / 31日 & サンジュウイチニチ \\
    \end{longtblr}

    \subsubsection{Les mois}

    Les mois de l'année sont formés à partir de la lecture « on » des douze premiers nombres japonais, suivis de ガツ.

    \begin{longtblr}{%
        colspec = {Q[l,2] Q[l,4] Q[l,6]},  
        width = \linewidth,
        hlines = {0.6pt, gray!50},
        vlines = {0pt},
        row{1} = {font=\bfseries, bg=couleurmatiere!20}
    }
        Mois & 漢字 & 仮名 \\
        Janvier & 一月 / 1月 & イチガツ \\
        Février & 二月 / 2月 & 二ガツ \\
        Mars & 三月 / 3月 & サンガツ \\
        Avril & 四月 / 4月 & シガツ \\
        Mai & 五月 / 5月 & ゴガツ \\
        Juin & 六月 / 6月 & ロクガツ \\
        Juillet & 七月 / 7月 & シチガツ \\
        Août & 八月 / 8月 & ハチガツ \\
        Septembre & 九月 / 9月 & クガツ \\
        Octobre & 十月 / 10月 & ジュウガツ \\
        Novembre & 十一月 / 11月 & ジュウイチガツ \\
        Décembre & 十二月 / 12月 & ジュウニガツ \\
    \end{longtblr}

    \begin{tcolorbox}
        -- 二千二十五年三月十五日

        -- 2025年3月15日

        -- ニセンニジュウゴネン　サンガツ　ジュウゴニチ

        -- 15/03/2025
    \end{tcolorbox}

    